\section{Controlled first crossing inverses}\label{sec:inverse}
For any of the four counting sequences, define
\[
 \nu_u(X)=\min\{n\ge0:u_n\ge X\}\qquad(X>1).
\]
Fresh-record padding gives an injection at every length and preserves each class, so every $u_n$ is nondecreasing. The lower bounds above make each sequence unbounded. Thus its first crossing exists; strict increase is unnecessary.

\begin{theorem}\label{thm:inverse}
Put $L=\log X$, $\lambda=\log L$, $\mu=\log\lambda$, and
\[
 D=\frac{L\mu^2}{\lambda^3}.
\]
For $u=a$ or $u=e$,
\begin{equation}\label{eq:inverseexact}
 \nu_u(X)=\frac L\lambda
       +(3\mu+1-\log2)\frac L{\lambda^2}+O(D).
\end{equation}
For $u=d$ or $u=c$, there is a fixed $C>0$, depending on the class, such that eventually
\begin{align}
 \frac L\lambda+(3\mu+1-\log2)\frac L{\lambda^2}-CD
 &\le\nu_u(X)\notag\\
 &\le\frac L\lambda+(3\mu+1+\log2)\frac L{\lambda^2}+CD.
 \label{eq:inversebracket}
\end{align}
These are asymptotic error bounds, not exact integer-rounding prescriptions.
\end{theorem}
\begin{proof}
For a fixed real constant $c$, define
\[
 f_c(x)=x(\log x-2\log\log x+c),\qquad
 x_c=\frac L\lambda\left(1+\frac{3\mu-c}{\lambda}\right).
\]
For sufficiently large $X$, all quantities are positive. Taylor expansion of the ordinary logarithm gives
\begin{align}
 \log x_c&=\lambda-\mu+\frac{3\mu-c}{\lambda}
                +O(\mu^2/\lambda^2),\notag\\
 \log\log x_c&=\mu-\frac\mu\lambda+O(\mu^2/\lambda^2),\notag\\
 f_c(x_c)&=L+O(L\mu^2/\lambda^2).\label{eq:inverseexpand}
\end{align}
For the middle line, divide the first line by $\lambda$ before expanding the second logarithm. Substitution then proves the last line: the order-$\mu/\lambda$ terms cancel by the choice $3\mu-c$.

On every fixed-width multiple of the window $x=x_c+O(D)$, one has $x\sim L/\lambda$ and
\begin{equation}\label{eq:slope}
 f_c'(x)=\log x-2\log\log x+c+1-2/\log x\sim\lambda.
\end{equation}
The logarithmic count error from Theorem~\ref{thm:main}, evaluated at any integer in that window, is
\begin{equation}\label{eq:inverseerror}
 O\!\left(\frac{x(\log\log x)^2}{\log x}\right)
       =O(L\mu^2/\lambda^2).
\end{equation}
The constants can be taken uniform there, since the original all-integer error bound holds for every sufficiently large integer and the displayed scales are uniformly comparable.

For $u=a,e$, take $c=C_+$. The mean-value theorem, \eqref{eq:inverseexpand}--\eqref{eq:inverseerror}, and a sufficiently large fixed $C$ imply
\[
 \log u_{\lfloor x_{C_+}-CD\rfloor}<L,
 \qquad
 \log u_{\lceil x_{C_+}+CD\rceil}\ge L.
\]
Rounding changes $f_c$ by only $O(\lambda)$, negligible compared with $L\mu^2/\lambda^2$. Monotonicity brackets the first crossing between these two neighboring integer endpoints. Because $D\to\infty$, their unit rounding errors are absorbed in $O(D)$. This proves \eqref{eq:inverseexact}.

For $u=d,c$, apply the upper count estimate with $c=C_+$ to the lower endpoint: it implies the count is still below $X$ at $\lfloor x_{C_+}-CD\rfloor$. Apply the lower count estimate with $c=C_-$ to the upper endpoint: the count is at least $X$ at $\lceil x_{C_-}+CD\rceil$. Choose one $C$ large enough for both estimates and absorb the unit rounding errors as before. This proves \eqref{eq:inversebracket}. In particular the directions reverse: the larger counting constant gives the smaller inverse endpoint.
\end{proof}
The error $D=L(\log\log L)^2/(\log L)^3$ tends to infinity in absolute size. Even for $a,e$, rounding the two displayed terms does not determine the true first crossing. The error is smaller than their second term, but is not bounded. In \eqref{eq:inversebracket}, the unresolved counting constant leaves an inverse interval of width
\[
 2\log2\,\frac L{\lambda^2}+O(D),
\]
and $D=o(L/\lambda^2)$. Thus the second coefficient for $d,c$ is genuinely undetermined by this argument. The constants implicit in the error terms and their eventual thresholds are not supplied as numerical certification tolerances.
